\documentclass{article}

\usepackage{amsmath}
\usepackage{amsthm}
\usepackage{amssymb}
\usepackage[hidelinks]{hyperref}

\newtheorem{lemma}{Lemma}
\newtheorem{theorem}{Theorem}
\providecommand*{\lemmaautorefname}{Lemma}

\newtheorem{corollary}[lemma]{Corollary}
\providecommand*{\corollaryautorefname}{Corollary}

\renewcommand{\refname}{\raggedright References}

\begin{document}

\begin{center}
\section*{Derivation of a Quantile Confidence Interval Based on Order Statistics}
{\small Frederik Danielsen}
\end{center}

\vspace{1em}

\noindent This note provides a mathematical derivation of a confidence interval for a quantile $\xi_q$ based on order statistics. The results presented here are classical and are not original to this note. They are derived independently to ensure a thorough understanding of the underlying mathematics, to present the arguments in a form suited to the Problab\footnote{\url{https://github.com/FrederikDanielsenDevelopment/problab}} project, and to obtain conclusions that are directly useful for the implementation in Problab. The derivation assumes i.i.d. samples $X_1,\ldots,X_n$ satisfying $P(X_i\leq\xi_q)=q$ and a significance level $\alpha\in(0,1)$, corresponding to a target coverage level of $1-\alpha$. In the first lemma, a confidence interval of the form $[X_{(k_L^*)},X_{(k_U^*+1)}]$ is derived by choosing the indices $k_L^*$ and $k_U^*$ such that the two binomial tail probabilities are bounded by $\alpha/2$. In the second lemma, it is shown that the iterative procedures used in the implementation terminate at exactly these optimal indices.

Practically, these results enable Problab to return finite-sample confidence intervals with coverage probability at least $1-\alpha$ for estimated quantiles, such as medians or percentiles.

\vspace{2em}

\begin{lemma}
Let $X_i$ with $i\in\{1,\ldots,n-1\}$ be i.i.d samples. Suppose that $P(X_i\leq \xi_q)=q$ for all $i$, some $q\in(0,1)$, and $\xi_q$ in the value space of $X_i$. Let
$$I_i =
\begin{cases}
1,\quad X_i\leq\xi_q,\\
0,\quad X_i>\xi_q.
\end{cases}$$
and let
$$K=\sum_{i=1}^n I_i\sim Bin(n,q).$$
Let $k_L^*$ and $k_U^*$ solve
$$
\begin{aligned}
k_L^*
&=
\max_{k_L\in\{1,\ldots,n-1\}}
\left\{
k_L:
P(K<k_L)\leq\frac{\alpha}{2}
\right\}\\
k_U^*
&=
\min_{k_U\in\{1,\ldots,n-1\}}
\left\{
k_U:
P(K>k_U)\leq\frac{\alpha}{2}
\right\}.
\end{aligned}
$$
Then
$$[X_{(k_L^*)},X_{(k_U^*+1)}]$$
is a confidence interval for $\xi_q$ with coverage probability at least $1-\alpha$. Furthermore, among the intervals satisfying the two tail constraints, solving the optimization problems above produces the narrowest interval.
\end{lemma}
\begin{proof}
$I_i$ can be regarded as the outcomes of individual Bernouli trials with probability of success $q$ with
$$K=\sum_{i=1}^nI_i\sim Bin(n,q)$$
denoting the number of samples with values at most $\xi_q$. Now, consider the ordered samples
$$X_{(1)},X_{(2)},\ldots,X_{(n)}.$$
If $K\leq k_U$ for some $k_U\in\{1,\ldots,n-1\}$, that would imply that at most $k_U$ samples fell below $\xi_q$. That is, $\xi_q<X_{(k_U+1)}$. Likewise, if $K\geq k_L$ for some $k_L\in\{1,\ldots,n-1\}$, that would imply that at least $k_L$ samples fell below $\xi_q$. That is, $\xi_q\geq X_{(k_L)}$. Hence,
$$P(k_L\leq K\leq k_U)=P(X_{(k_L)}\leq\xi_q<X_{(k_U+1)}).$$
Therefore, if we determine $k_L$ and $k_U$ such that
$$P(k_L\leq K \leq k_U)=1-\alpha,$$
for some $\alpha\in(0,1)$, then the interval $[X_{(k_L)},X_{(k_U+1)}]$ is a confidence interval for $\xi_q$ with coverage probability $1-\alpha$. However, this equality is one with two unknowns and thus multiple solutions. However, we may impose an additional constraint, namely that the likelihood that $\xi_q$ lies above the confidence interval is the same as the likelihood that it lies below it $-$ at least to the extent that a discrete number of samples allows it. Specifically, we require that
$$P(\xi_q\geq X_{(k_U+1)})=P(\xi_q<X_{(k_L)})=\alpha/2.$$
However, as mentioned, since we are dealing with a discrete number of samples, it is likely that no $(k_L,k_U)$ exists that fulfil this exactly. Therefore, we make do with the best estimate (likely more conservative than $1-\alpha$) and require simply that
$$
\begin{aligned}
P(\xi_q \geq X_{(k_U+1)})&\leq \alpha/2,\\
P(\xi_q<X_{(k_L)})&\leq \alpha/2\\
\Leftrightarrow\\
P(K>k_U)&\leq \alpha/2,\\
P(K<k_L)&\leq \alpha/2.
\end{aligned}
$$
Since $\alpha<1$, any pair $(k_L,k_U)$ satisfying both tail constraints necessarily satisfies $k_L\leq k_U$. Indeed, if $k_L>k_U$, then every possible value of $K$ satisfies either $K<k_L$ or $K>k_U$, and therefore
$$
1=P\big((K<k_L)\cup(K>k_U)\big)
\leq P(K<k_L)+P(K>k_U)
\leq \alpha,
$$
which is a contradiction. Furthermore, in case you were wondering, since the lower endpoint $X_{(k_L)}$ requires $k_L\geq1$ and the upper endpoint $X_{(k_U+1)}$ requires $k_U\leq n-1$, the inequality $k_L\leq k_U$ implies that any feasible pair must satisfy
$$
k_L,k_U\in\{1,\ldots,n-1\}.
$$
To obtain the narrowest confidence interval satisfying these constraints, we choose the largest possible $k_L$ and the smallest possible $k_U$. Thus, we solve
$$
\begin{aligned}
k_L^*
&=
\max_{k_L\in\{1,\ldots,n-1\}}
\left\{
k_L:
P(K<k_L)\leq\frac{\alpha}{2}
\right\},\\
k_U^*
&=
\min_{k_U\in\{1,\ldots,n-1\}}
\left\{
k_U:
P(K>k_U)\leq\frac{\alpha}{2}
\right\}.
\end{aligned}
$$
In this way we ensure that
$$P(k_L^*\leq K \leq k_U^*)=P(X_{(k_L^*)}\leq \xi_q< X_{(k_U^*+1)})\geq 1-\alpha.$$
Then the interval
$$[X_{(k_L^*)},X_{(k_U^*+1)}]$$
is at least as conservative as a confidence interval with coverage probability $1-\alpha$.
\end{proof}

\begin{lemma}
Let $K\sim Bin(n,q)$ and $\alpha\in(0,1)$, and suppose that the optimization problems
$$
\begin{aligned}
k_L^*
&=
\max_{k\in\{1,\ldots,n-1\}}
\left\{
k:
P(K<k)\leq\frac{\alpha}{2}
\right\},\\
k_U^*
&=
\min_{k\in\{1,\ldots,n-1\}}
\left\{
k:
P(K>k)\leq\frac{\alpha}{2}
\right\}
\end{aligned}
$$
have solutions. Consider the iterative procedures
$$
\begin{aligned}
&k_L\leftarrow 1,\\
&\text{while } k_L<n-1
\text{ and }
P(K\leq k_L)\leq\frac{\alpha}{2},
\quad
k_L\leftarrow k_L+1,
\end{aligned}
$$
and
$$
\begin{aligned}
&k_U\leftarrow n-1,\\
&\text{while } k_U>1
\text{ and }
P(K\geq k_U)\leq\frac{\alpha}{2},
\quad
k_U\leftarrow k_U-1.
\end{aligned}
$$
Then the procedures terminate with
$$
k_L=k_L^*
\qquad\text{and}\qquad
k_U=k_U^*.
$$
\end{lemma}
\begin{proof}
For the lower index, observe that increasing the current value $k_L$ to $k_L+1$ is permissible exactly when
$$
P(K<k_L+1)=P(K\leq k_L)\leq\frac{\alpha}{2}.
$$
Thus, at every iteration, the algorithm increments $k_L$ only when the next larger value satisfies the lower-tail constraint. Since $k_L$ increases by one at each iteration and is bounded above by $n-1$, the procedure terminates.
If the procedure stops before reaching $n-1$, then increasing $k_L$ by one would violate the constraint. Indeed, the next candidate would be $k_L+1$, and for this candidate the constraint is
$$
P(K<k_L+1)\leq\frac{\alpha}{2}.
$$
Since
$$
P(K<k_L+1)=P(K\leq k_L),
$$
and the loop has stopped, this inequality is no longer satisfied. Therefore, $k_L+1$ is not feasible. Since $P(K<k)$ is nondecreasing as $k$ increases, no larger value can be feasible either. Hence, the returned $k_L$ is the largest feasible value, so
$$
k_L
=
\max_{k\in\{1,\ldots,n-1\}}
\left\{
k:
P(K<k)\leq\frac{\alpha}{2}
\right\}
=
k_L^*.
$$
If the procedure instead terminates at $k_L=n-1$, then $n-1$ is itself the largest possible value in the domain, so again $k_L=k_L^*$.
For the upper index, observe that decreasing the current value $k_U$ to $k_U-1$ is permissible exactly when
$$
P(K>k_U-1)=P(K\geq k_U)\leq\frac{\alpha}{2}.
$$
Thus, at every iteration, the algorithm decrements $k_U$ only when the next smaller value satisfies the upper-tail constraint. Since $k_U$ decreases by one at each iteration and is bounded below by $1$, the procedure terminates.
If the procedure stops before reaching $1$, then decreasing $k_U$ by one would violate the constraint. Indeed, the next candidate would be $k_U-1$, and for this candidate the constraint is
$$
P(K>k_U-1)\leq\frac{\alpha}{2}.
$$
Since
$$
P(K>k_U-1)=P(K\geq k_U),
$$
and the loop has stopped, this inequality is no longer satisfied. Therefore, $k_U-1$ is not feasible. Since $P(K>k)$ increases as $k$ decreases, no smaller value can be feasible either. Hence, the returned $k_U$ is the smallest feasible value, so
$$
k_U
=
\min_{k\in\{1,\ldots,n-1\}}
\left\{
k:
P(K>k)\leq\frac{\alpha}{2}
\right\}
=
k_U^*.
$$
If the procedure instead terminates at $k_U=1$, then $1$ is itself the smallest possible value in the domain, so again $k_U=k_U^*$.
\end{proof}



\end{document}

